% MT058 / explanation — canonical editable content.
Bernoulli NB把一句话变成一串0/1：词表中的词出现过记1，没出现记0。“happy”出现一次和三次，在这项特征上没有区别。词袋还忽略顺序，因此“happy”和“not happy”的组合含义不容易表达。改进要留在算力有限的词袋生成式路线内。


\textbf{第一项：尝试词计数与Multinomial NB。} 改用每个词的出现次数，使重复词能够多次提供证据。训练估计各情绪类别中的词概率，预测时仍主要是查表和加法。重复未必总有用，所以用验证数据比较。


\textbf{第二项：加入少量有意义的bigram。} Bigram是相邻两个词组成的特征，例如把“not happy”作为一个整体计数。这补充了局部顺序，代价是词表更大；在设备约束下筛选常见或有效的短语，不能无限扩充。


\textbf{第三项：调平滑强度。} 训练只有5000句，某词在某类未出现并不表示它永远不可能出现。加性平滑给计数加一个小量，使概率不轻易为0。平滑太强又会把不同类别的词分布抹得太像，因此强度需要验证。


\textbf{第四项：选择合适词表。} 去除几乎没有信息或极不可靠的词，保留情绪相关词，并比较不同词表规模。这同时控制内存、推理成本和估计噪声。不要自动删掉not，否定词可能改变整句含义；感叹号也可能提供情绪信息。


词表与概率由训练集确定，用验证集选择上述方案，最后评价测试集。


\textbf{出处：} 2023B期中 Q7，原卷第4页（10分）


\textbf{答案依据：} 2023B期中Q7随卷解答；不把更复杂模型当必然更好。
